% This file is included by uet_thesis.tex.
% Replace blue placeholders and remove example text before submission.

\chapter{Design Specification}
\label{chap:design}
% =============================================================================

\section{System Context, Architecture, and Interfaces \applicability{Development, Hybrid, Game; Research Software Prototype}}
\begin{guidance}
Define the system boundary, external actors or systems, major components, their
responsibilities, and principal interfaces. State the selected architectural style
only when it helps explain the design. Label every component and connector.
\end{guidance}

\diagram[dy]{system-context}{System context diagram}
\diagram[dy]{high-level-architecture}{High-level architecture diagram}

\begin{xltabular}{\textwidth}{@{}p{0.13\textwidth}p{0.24\textwidth}X X@{}}
\caption{Major interfaces}\label{tab:major-interfaces}\\
\toprule
\textbf{ID} & \textbf{Interface} & \textbf{Purpose and exchanged information} & \textbf{Protocol or format} \\
\midrule
\endfirsthead
\multicolumn{4}{l}{\small\textit{(Continued from previous page)}}\\
\toprule
\textbf{ID} & \textbf{Interface} & \textbf{Purpose and exchanged information} & \textbf{Protocol or format} \\
\midrule
\endhead
\multicolumn{4}{r}{\small\textit{(Continued on next page)}}\\
\endfoot
\bottomrule
\endlastfoot
INT-01 & CLI $\leftrightarrow$ Interpreter & Invoke parse, compile, validate, and extract-source on \texttt{.md} files and templates. & Python API/local calls \\
INT-02 & Interpreter $\leftrightarrow$ Tectonic & Pass generated \texttt{.tex}, receive compiled PDF; supply template-required packages. & Process invocation / filesystem \\
INT-03 & Editor Extensions $\leftrightarrow$ Template Schema/Manifest & Load grammar, snippets, and IntelliSense from template schema and manifest. & Local files / extension APIs \\
INT-04 & Marketplace $\leftrightarrow$ GitHub & Read labeled issues for the published template registry; resolve template manifests and assets from linked repositories. & GitHub REST API \\
INT-05 & MCP Server $\leftrightarrow$ AI Assistants & Expose build, validate, template retrieval/generation/modification, specification-to-Markdown conversion, and extract-source operations guided by \texttt{skills.md}. & MCP (JSON-RPC) \\
INT-06 & Conversion API $\leftrightarrow$ CMS/Websites & Accept Markdown + template, return PDF or error; support webhooks. & HTTP/REST (JSON) \\
INT-07 & PDF $\leftrightarrow$ Extract-Source Tool & Read/write embedded Markdown source in PDF metadata. & PDF metadata (XMP) \\
INT-08 & Conversion Web App $\leftrightarrow$ Interpreter/API & Submit Markdown with a selected template; return a compiled PDF or restored source. & HTTPS/JSON \\
INT-09 & Documentation Website $\leftrightarrow$ Repository & Publish versioned installation, configuration, and usage guides from the documentation source. & Static build / Git \\
INT-10 & Marketplace $\leftrightarrow$ Counters Database & Read and increment per-template heart and download counts. & SQL/ORM \\
\end{xltabular}


\section{Behavioural and Interaction Design \applicability{Development, Hybrid, Game; Research if applicable}}
\begin{guidance}
Include only the models needed to explain the project's behaviour and interaction.
Each UML diagram must use standard UML notation~\citep{omguml251}, remain readable, and be followed by a
brief explanation of important flows, alternatives, or design decisions.
\end{guidance}

\subsection{Use-Case Model}
\begin{guidance}
Show external actors, system use cases, the system boundary, and justified
\texttt{include}, \texttt{extend}, or generalization relationships. Use cases
should represent user goals rather than interface screens.
\end{guidance}
\diagram[dy]{use-case}{System-wide use-case diagram}

The Litho publishing ecosystem encompasses seventeen formal use cases organized across five core functional subsystems, supporting five distinct external actors:
\begin{enumerate}
    \item \textbf{Author / Student:} Technical writers and researchers creating documents, compiling PDFs, inspecting marketplace templates, and extracting source from compiled artifacts.
    \item \textbf{Template Publisher:} Institutional formatting leads and designers authoring, validating, and submitting reproducible template packages.
    \item \textbf{Marketplace Maintainer:} Curators who review automated security baselines, attest capabilities, and publish approved listings.
    \item \textbf{AI Assistant:} Autonomous agents interacting with Litho tools over the Model Context Protocol (MCP).
    \item \textbf{CMS Site / External Client:} Third-party content platforms and integration pipelines invoking automated conversion over REST APIs and webhooks.
\end{enumerate}

Each use-case specification below is structured according to its key architectural components: Actor, System, Goal (Objective), Preconditions, Trigger, Basic Flow (Main Scenario), and Alternate / Exception Flows, as presented in Tables~\ref{tab:uc-01} through~\ref{tab:uc-17}.

\begin{xltabular}{\textwidth}{@{}p{0.26\textwidth}X@{}}
\caption{Use-case specification: UC-01 Build PDF Document Locally}\label{tab:uc01}\\
\toprule
\textbf{Component} & \textbf{Description} \\
\midrule
\endfirsthead
\multicolumn{2}{l}{\small\textit{(Continued from previous page)}}\\
\toprule
\textbf{Component} & \textbf{Description} \\
\midrule
\endhead
\bottomrule
\endlastfoot
\textbf{Actor} & Author / Student (Secondary: Tectonic Typesetting Engine) \\
\textbf{System} & Litho CLI and Typesetting Pipeline \\
\textbf{Goal (Objective)} & Compile a Litho Markdown document and its bound template into a publication-quality PDF containing embedded source and provenance metadata. \\
\textbf{Preconditions} & Markdown source file and referenced template exist locally; Litho CLI or editor extension is installed and configured. \\
\textbf{Trigger} & Author executes \texttt{litho build} command or triggers compilation from editor. \\
\textbf{Basic Flow (Main Scenario)} & 1.~Author executes \texttt{litho build} command targeting document file; 2.~Tokenizer and lexer read source file and construct the three-tier Abstract Syntax Tree (AST); 3.~Schema validator inspects document frontmatter and AST structure against template constraints; 4.~Template compiler binds AST content to target \LaTeX{} macros, producing \texttt{.tex} source; 5.~Tectonic engine compiles generated \texttt{.tex} into PDF binary; 6.~Metadata engine embeds raw Markdown source, template ID, and version into PDF XMP stream; 7.~System writes compiled PDF to disk and displays build summary with page count and execution time. \\
\textbf{Alternate / Exception Flows} & \textbf{2a.~Syntax Error:} Lexer encounters invalid Markdown dialect or malformed blocks; halts build and reports line/column diagnostics. \textbf{3a.~Schema Failure:} Frontmatter violates required template fields; halts execution and displays constraint violations. \textbf{5a.~\LaTeX{} Error:} Tectonic encounters typesetting error; parses error log and returns structured diagnostics without overwriting previous PDF. \\
\textbf{Related Requirements} & FR-01--FR-08, FR-10, FR-51, FR-52 \\
\end{xltabular}

\begin{xltabular}{\textwidth}{@{}p{0.26\textwidth}X@{}}
\caption{Use-case specification: UC-02 Edit Markdown with Real-Time Diagnostics}\label{tab:uc02}\\
\toprule
\textbf{Component} & \textbf{Description} \\
\midrule
\endfirsthead
\multicolumn{2}{l}{\small\textit{(Continued from previous page)}}\\
\toprule
\textbf{Component} & \textbf{Description} \\
\midrule
\endhead
\bottomrule
\endlastfoot
\textbf{Actor} & Author / Student (Secondary: VS Code Extension Language Service) \\
\textbf{System} & Litho VS Code Extension \& Language Service \\
\textbf{Goal (Objective)} & Author Litho Markdown documents with syntax highlighting, template shortcode completion, and live inline error diagnostics. \\
\textbf{Preconditions} & VS Code is installed with the Litho extension enabled; a Litho Markdown document is open in the workspace. \\
\textbf{Trigger} & Author opens, types, or edits content in a \texttt{.md} file. \\
\textbf{Basic Flow (Main Scenario)} & 1.~Author opens or edits document in VS Code; 2.~Extension recognizes Litho grammar and applies contextual syntax highlighting; 3.~Language service resolves referenced template manifest from document frontmatter; 4.~Language service provides snippet insertion and IntelliSense auto-completion for template-defined shortcodes and parameters; 5.~Incremental background parser continuously validates AST syntax and frontmatter schemas; 6.~Language service publishes inline warning and error squiggles to the editor Problems panel. \\
\textbf{Alternate / Exception Flows} & \textbf{3a.~Missing Manifest:} Extension displays warning notification, falls back to core Litho syntax highlighting, and disables template-specific shortcode completion. \textbf{5a.~Schema Violation:} Problem marker highlights invalid frontmatter field with expected type and allowed values. \\
\textbf{Related Requirements} & FR-16, FR-17, FR-18, FR-19 \\
\end{xltabular}

\begin{xltabular}{\textwidth}{@{}p{0.26\textwidth}X@{}}
\caption{Use-case specification: UC-03 Watch and Hot-Rebuild Document}\label{tab:uc03}\\
\toprule
\textbf{Component} & \textbf{Description} \\
\midrule
\endfirsthead
\multicolumn{2}{l}{\small\textit{(Continued from previous page)}}\\
\toprule
\textbf{Component} & \textbf{Description} \\
\midrule
\endhead
\bottomrule
\endlastfoot
\textbf{Actor} & Author / Student (Secondary: OS File Watcher Daemon) \\
\textbf{System} & Litho CLI Watch Daemon \& Fast Rebuild Engine \\
\textbf{Goal (Objective)} & Automatically recompile and refresh output PDF upon saving changes to source Markdown, assets, or template files. \\
\textbf{Preconditions} & Initialized Litho project directory containing document source and template files. \\
\textbf{Trigger} & Author executes \texttt{litho watch} command in the project directory. \\
\textbf{Basic Flow (Main Scenario)} & 1.~Author executes \texttt{litho watch}; 2.~System performs initial compilation and registers OS-level file system change watchers on \texttt{.md}, asset, and template files; 3.~Author edits and saves a project file; 4.~Watch daemon receives file modification event and triggers incremental re-compilation; 5.~System re-validates modified AST, regenerates \LaTeX{}, and recompiles PDF via Tectonic; 6.~System updates PDF binary on disk and prints rebuild timestamp and duration to console. \\
\textbf{Alternate / Exception Flows} & \textbf{4a.~Incremental Build Error:} Compiler detects syntax or \LaTeX{} compilation failure; logs diagnostics to console, keeps prior valid PDF on disk, and remains active watching for file fixes. \\
\textbf{Related Requirements} & FR-11 \\
\end{xltabular}

\begin{xltabular}{\textwidth}{@{}p{0.26\textwidth}X@{}}
\caption{Use-case specification: UC-04 Restore Markdown Source from PDF}\label{tab:uc04}\\
\toprule
\textbf{Component} & \textbf{Description} \\
\midrule
\endfirsthead
\multicolumn{2}{l}{\small\textit{(Continued from previous page)}}\\
\toprule
\textbf{Component} & \textbf{Description} \\
\midrule
\endhead
\bottomrule
\endlastfoot
\textbf{Actor} & Author / Student (Secondary: PDF Metadata Stream Extractor) \\
\textbf{System} & Litho CLI \& Extract-Source Utility \\
\textbf{Goal (Objective)} & Extract and restore the exact original, byte-for-byte Markdown source code and template provenance from a compiled Litho PDF. \\
\textbf{Preconditions} & A Litho-generated PDF file exists and is accessible to the system. \\
\textbf{Trigger} & Author executes \texttt{litho extract <filename.pdf>} command. \\
\textbf{Basic Flow (Main Scenario)} & 1.~Author invokes extract command providing path to a compiled PDF; 2.~System parses the PDF document catalog and inspects embedded XMP metadata and attachment streams; 3.~System verifies Litho provenance signature, template identifier, and template version; 4.~System extracts embedded Markdown stream, verifies payload integrity hash, and decompresses content; 5.~System writes recovered \texttt{.md} file to disk and reports provenance metadata (template name, version, and original build timestamp). \\
\textbf{Alternate / Exception Flows} & \textbf{2a.~Missing Provenance Signature:} PDF lacks Litho metadata streams; system halts and displays error indicating file was not built with Litho or metadata was stripped. \textbf{4a.~Integrity Hash Mismatch:} Extracted source payload fails checksum verification; system alerts user of potential tampering or file corruption. \\
\textbf{Related Requirements} & FR-06, FR-07, FR-08, FR-12, FR-27, NFR-05 \\
\end{xltabular}

\begin{xltabular}{\textwidth}{@{}p{0.26\textwidth}X@{}}
\caption{Use-case specification: UC-05 Browse and Inspect Marketplace Templates}\label{tab:uc05}\\
\toprule
\textbf{Component} & \textbf{Description} \\
\midrule
\endfirsthead
\multicolumn{2}{l}{\small\textit{(Continued from previous page)}}\\
\toprule
\textbf{Component} & \textbf{Description} \\
\midrule
\endhead
\bottomrule
\endlastfoot
\textbf{Actor} & Author / Student (Secondary: Marketplace Static Catalog) \\
\textbf{System} & Litho Template Marketplace Web Client \\
\textbf{Goal (Objective)} & Discover, search, filter, and inspect curated, verified Litho document templates, viewing visual previews, metadata, and install instructions. \\
\textbf{Preconditions} & User has network access to the Litho Template Marketplace web application. \\
\textbf{Trigger} & User navigates to the marketplace homepage or catalog explorer. \\
\textbf{Basic Flow (Main Scenario)} & 1.~User navigates to marketplace catalog interface; 2.~Web client fetches static \texttt{catalog.json} and renders indexed template cards with titles, tags, and thumbnails; 3.~User enters search keywords or filters templates by category (e.g., Thesis, IEEE Paper, Book); 4.~Client applies client-side search filtering across template titles, descriptions, and tags; 5.~User clicks a template card to view detailed inspection page; 6.~System renders template details: author provenance, semantic version, license, verified badge, sample PDF preview, and CLI install command. \\
\textbf{Alternate / Exception Flows} & \textbf{4a.~No Search Matches:} Client displays empty search state with clear-filter suggestion. \textbf{6a.~Preview Asset Failure:} Client displays fallback placeholder and technical metadata summary. \\
\textbf{Related Requirements} & FR-20, FR-21, FR-23, NFR-10 \\
\end{xltabular}

\begin{xltabular}{\textwidth}{@{}p{0.26\textwidth}X@{}}
\caption{Use-case specification: UC-06 React and Download Template}\label{tab:uc06}\\
\toprule
\textbf{Component} & \textbf{Description} \\
\midrule
\endfirsthead
\multicolumn{2}{l}{\small\textit{(Continued from previous page)}}\\
\toprule
\textbf{Component} & \textbf{Description} \\
\midrule
\endhead
\bottomrule
\endlastfoot
\textbf{Actor} & Author / Student (Secondary: Cloudflare Worker \& D1 Engagement Service) \\
\textbf{System} & Marketplace Engagement Client \& Edge Service \\
\textbf{Goal (Objective)} & Support a community template with an anonymous heart reaction and copy the standard installation command or download the template archive. \\
\textbf{Preconditions} & User is viewing a specific template detail page in the marketplace. \\
\textbf{Trigger} & User clicks the heart reaction button or the copy install command button. \\
\textbf{Basic Flow (Main Scenario)} & 1.~User clicks the heart button to react to the template; 2.~Browser verifies local storage state to guard against duplicate reactions and updates UI count optimistically; 3.~Client transmits anonymous signed POST request to Cloudflare Worker edge endpoint; 4.~Worker validates sliding-window rate limit and atomically increments template heart counter in D1 database; 5.~User clicks ``Copy Install Command'' button; 6.~Client copies \texttt{litho install <template-id>} to system clipboard, shows confirmation badge, and increments copy telemetry counter. \\
\textbf{Alternate / Exception Flows} & \textbf{2a.~Duplicate Reaction:} Client detects prior heart in local storage; toggles off reaction and decrements local count. \textbf{4a.~Rate Limit Exceeded:} Worker returns HTTP 429; client reverts optimistic UI update gracefully. \\
\textbf{Related Requirements} & FR-53, FR-54, FR-55, FR-56 \\
\end{xltabular}

\begin{xltabular}{\textwidth}{@{}p{0.26\textwidth}X@{}}
\caption{Use-case specification: UC-07 Convert Markdown in Web Application}\label{tab:uc07}\\
\toprule
\textbf{Component} & \textbf{Description} \\
\midrule
\endfirsthead
\multicolumn{2}{l}{\small\textit{(Continued from previous page)}}\\
\toprule
\textbf{Component} & \textbf{Description} \\
\midrule
\endhead
\bottomrule
\endlastfoot
\textbf{Actor} & Author / Student (Secondary: Hosted Typesetting Sandbox Service) \\
\textbf{System} & Litho Web Conversion Application \\
\textbf{Goal (Objective)} & Author or upload Markdown in the browser, choose a template, and generate a downloadable PDF without installing local CLI or \LaTeX{} tools. \\
\textbf{Preconditions} & Web browser with internet connectivity to the conversion web portal. \\
\textbf{Trigger} & User pastes Markdown text or uploads a \texttt{.md} file in the web editor and clicks ``Convert to PDF''. \\
\textbf{Basic Flow (Main Scenario)} & 1.~User inputs Markdown text into the web editor or uploads an existing document; 2.~User selects an official or community template from the dropdown selector; 3.~User clicks ``Convert to PDF'' button; 4.~Web client submits document payload and template selection to the sandboxed conversion backend; 5.~Backend executes three-tier AST parser, binds template, and compiles PDF within an isolated process sandbox; 6.~Backend embeds Markdown source into output PDF metadata; 7.~Backend streams generated PDF binary back to client; web client displays PDF in inline preview pane with download link. \\
\textbf{Alternate / Exception Flows} & \textbf{4a.~Input Size Exceeded:} Client flags size error before network dispatch. \textbf{5a.~Compilation Fails in Sandbox:} Sandbox captures compiler diagnostic log, maps error to source lines, and displays structured error banner in editor. \\
\textbf{Related Requirements} & FR-24, FR-25, FR-26, FR-27, NFR-06 \\
\end{xltabular}

\begin{xltabular}{\textwidth}{@{}p{0.26\textwidth}X@{}}
\caption{Use-case specification: UC-08 Convert Technical Specification via AI}\label{tab:uc08}\\
\toprule
\textbf{Component} & \textbf{Description} \\
\midrule
\endfirsthead
\multicolumn{2}{l}{\small\textit{(Continued from previous page)}}\\
\toprule
\textbf{Component} & \textbf{Description} \\
\midrule
\endhead
\bottomrule
\endlastfoot
\textbf{Actor} & Author / Student (Secondary: LLM Inference Service) \\
\textbf{System} & Litho AI Authoring Assistant \\
\textbf{Goal (Objective)} & Transform unstructured technical notes, SRS specifications, or outlines into a fully formatted Litho Markdown document tailored to a selected template. \\
\textbf{Preconditions} & User has opened the AI authoring modal in the web app and provided API key or session access. \\
\textbf{Trigger} & User inputs raw technical specification notes and clicks ``Generate Litho Document''. \\
\textbf{Basic Flow (Main Scenario)} & 1.~User enters raw specification text and selects target template (e.g., University Thesis or IEEE Paper); 2.~AI subsystem loads the template manifest schema and the domain-specific \texttt{skills.md} agent playbook into context; 3.~Subsystem crafts structured prompt instructing LLM to map technical entities into Litho frontmatter, sectioning, and shortcodes; 4.~LLM infers and generates clean Litho Markdown matching template constraints; 5.~System validates generated Markdown through AST parser; 6.~System populates editor pane with validated Markdown and presents preview to author. \\
\textbf{Alternate / Exception Flows} & \textbf{5a.~Invalid Shortcode Generated:} AST validator detects unmapped shortcode tags; prompts model to correct syntax using schema rules before displaying to user. \\
\textbf{Related Requirements} & FR-28, FR-29, FR-30, FR-32, FR-33, FR-34 \\
\end{xltabular}

\begin{xltabular}{\textwidth}{@{}p{0.26\textwidth}X@{}}
\caption{Use-case specification: UC-09 Create Template Package}\label{tab:uc09}\\
\toprule
\textbf{Component} & \textbf{Description} \\
\midrule
\endfirsthead
\multicolumn{2}{l}{\small\textit{(Continued from previous page)}}\\
\toprule
\textbf{Component} & \textbf{Description} \\
\midrule
\endhead
\bottomrule
\endlastfoot
\textbf{Actor} & Template Publisher \\
\textbf{System} & Litho CLI Template Scaffolding Tool \\
\textbf{Goal (Objective)} & Scaffold a new, compliant Litho template package containing \texttt{manifest.json}, \LaTeX{} macro definitions, declarative schema, and sample document. \\
\textbf{Preconditions} & Litho CLI installed; publisher is in target development workspace. \\
\textbf{Trigger} & Publisher runs \texttt{litho template init <template-name>} command. \\
\textbf{Basic Flow (Main Scenario)} & 1.~Publisher executes template initialization command; 2.~CLI prompts publisher for template title, description, category, license, and author details; 3.~System generates template directory structure: \texttt{manifest.json}, \texttt{template.tex}, \texttt{schema.json}, preview assets, and \texttt{example.md}; 4.~System populates \texttt{manifest.json} with semantic version \texttt{1.0.0} and default configuration; 5.~CLI confirms successful scaffolding and outputs recommended next steps for template customization. \\
\textbf{Alternate / Exception Flows} & \textbf{2a.~Directory Conflict:} System detects destination directory already exists; prompts publisher to choose a different name or force overwrite. \\
\textbf{Related Requirements} & FR-13, FR-44--FR-47, NFR-04 \\
\end{xltabular}

\begin{xltabular}{\textwidth}{@{}p{0.26\textwidth}X@{}}
\caption{Use-case specification: UC-10 Validate Template Package}\label{tab:uc10}\\
\toprule
\textbf{Component} & \textbf{Description} \\
\midrule
\endfirsthead
\multicolumn{2}{l}{\small\textit{(Continued from previous page)}}\\
\toprule
\textbf{Component} & \textbf{Description} \\
\midrule
\endhead
\bottomrule
\endlastfoot
\textbf{Actor} & Template Publisher \\
\textbf{System} & Litho Template Linter \& Static Validator \\
\textbf{Goal (Objective)} & Run automated verification checks on a local template package to ensure manifest completeness, schema validity, and safe \LaTeX{} macro usage. \\
\textbf{Preconditions} & Template files created and accessible in the publisher workspace. \\
\textbf{Trigger} & Publisher executes \texttt{litho template validate} command. \\
\textbf{Basic Flow (Main Scenario)} & 1.~Publisher executes validation command targeting template folder; 2.~Validator verifies \texttt{manifest.json} against canonical JSON schema (ID, version, author, shortcode definitions); 3.~Validator parses \texttt{template.tex} to confirm required placeholder insertion hooks exist; 4.~Validator runs security static analysis to ensure no dangerous \LaTeX{} primitives (e.g., \texttt{\textbackslash write18}, shell escapes, pipe operations) are present; 5.~Validator executes a test compilation with \texttt{example.md} to guarantee compilation success; 6.~CLI displays green validation pass badge and outputs compliance report. \\
\textbf{Alternate / Exception Flows} & \textbf{2a.~Manifest Schema Violation:} Validator flags missing mandatory fields or invalid semantic versioning format. \textbf{4a.~Forbidden Macro Detected:} Validator flags banned macro command with exact line number; compilation test is blocked until removed. \\
\textbf{Related Requirements} & FR-14, NFR-08, NFR-09 \\
\end{xltabular}

\begin{xltabular}{\textwidth}{@{}p{0.26\textwidth}X@{}}
\caption{Use-case specification: UC-11 Submit Template via GitHub Issue}\label{tab:uc11}\\
\toprule
\textbf{Component} & \textbf{Description} \\
\midrule
\endfirsthead
\multicolumn{2}{l}{\small\textit{(Continued from previous page)}}\\
\toprule
\textbf{Component} & \textbf{Description} \\
\midrule
\endhead
\bottomrule
\endlastfoot
\textbf{Actor} & Template Publisher (Secondary: GitHub REST API \& Actions) \\
\textbf{System} & Litho Submission Gateway \& GitHub Workflow \\
\textbf{Goal (Objective)} & Submit a verified, public template Git repository to the Litho Template Marketplace registry for automated scanning and listing review. \\
\textbf{Preconditions} & Template repository is public on GitHub, passes local validation, and has tagged release. \\
\textbf{Trigger} & Publisher submits a structured GitHub Issue using the official template submission form. \\
\textbf{Basic Flow (Main Scenario)} & 1.~Publisher opens new submission Issue on marketplace repository; 2.~Publisher fills in repository URL, commit SHA, release tag, category, and confirmation checklist; 3.~GitHub Actions submission listener webhook triggers upon issue creation; 4.~Automated runner clones upstream repository at exact specified commit SHA; 5.~Runner executes Automated Security Baseline scanner and structural validation suite; 6.~Workflow posts bot report comment on issue indicating outcome (\texttt{passed}, \texttt{review-required}, or \texttt{needs-fixes}). \\
\textbf{Alternate / Exception Flows} & \textbf{2a.~Malformed Issue Form:} Bot flags missing mandatory fields in issue body and requests publisher update. \textbf{5a.~Security Baseline Finding Detected:} Bot comments detailed diagnostic finding log on issue and marks status as \texttt{needs-fixes} or \texttt{review-required}. \\
\textbf{Related Requirements} & FR-15, FR-22 \\
\end{xltabular}

\begin{xltabular}{\textwidth}{@{}p{0.26\textwidth}X@{}}
\caption{Use-case specification: UC-12 Generate Template via AI Assistance}\label{tab:uc12}\\
\toprule
\textbf{Component} & \textbf{Description} \\
\midrule
\endfirsthead
\multicolumn{2}{l}{\small\textit{(Continued from previous page)}}\\
\toprule
\textbf{Component} & \textbf{Description} \\
\midrule
\endhead
\bottomrule
\endlastfoot
\textbf{Actor} & Template Publisher (Secondary: LLM Synthesis Engine) \\
\textbf{System} & Litho AI Template Engineering Studio \\
\textbf{Goal (Objective)} & Rapidly synthesize a new template package from natural language style requirements or target PDF visual samples. \\
\textbf{Preconditions} & Publisher is authenticated in AI studio with prompt access enabled. \\
\textbf{Trigger} & Publisher provides layout guidelines (e.g., margins, typography, chapter header styles) and prompts AI to generate a template. \\
\textbf{Basic Flow (Main Scenario)} & 1.~Publisher inputs style specifications and target publication guidelines; 2.~AI engine retrieves canonical template blueprints and safe \LaTeX{} macro primitives; 3.~Engine synthesizes \texttt{manifest.json}, \texttt{template.tex}, and \texttt{schema.json} compliant with Litho architecture; 4.~System automatically runs validation checks on the generated package; 5.~System compiles sample test document and renders visual preview of output pages; 6.~Publisher inspects rendered layout and exports generated template package to repository. \\
\textbf{Alternate / Exception Flows} & \textbf{4a.~Synthetic Code Fails Validation:} Engine automatically feeds validation errors back into LLM prompt loop for automated iterative repair up to 3 attempts. \\
\textbf{Related Requirements} & FR-28, FR-31, FR-33, FR-34 \\
\end{xltabular}

\begin{xltabular}{\textwidth}{@{}p{0.26\textwidth}X@{}}
\caption{Use-case specification: UC-13 Perform Automated Security Review}\label{tab:uc13}\\
\toprule
\textbf{Component} & \textbf{Description} \\
\midrule
\endfirsthead
\multicolumn{2}{l}{\small\textit{(Continued from previous page)}}\\
\toprule
\textbf{Component} & \textbf{Description} \\
\midrule
\endhead
\bottomrule
\endlastfoot
\textbf{Actor} & Maintainer / Automated CI Scanner \\
\textbf{System} & Litho Automated Security Baseline Scanner \\
\textbf{Goal (Objective)} & Statically scan submitted template source repositories at the exact validated commit without executing arbitrary code, enforcing deterministic security outcomes. \\
\textbf{Preconditions} & Triggered by submission issue, update request, or manual audit dispatch. \\
\textbf{Trigger} & New submission issue created or maintainer applies \texttt{/scan} command. \\
\textbf{Basic Flow (Main Scenario)} & 1.~Scanner resolves submitted GitHub repository URL and binds exact commit SHA; 2.~Scanner inspects directory tree structure and manifest integrity; 3.~Static analyzer parses all \texttt{.tex}, \texttt{.sty}, and configuration scripts against forbidden execution signatures (shell escape, binary payloads, external script fetches); 4.~Scanner computes deterministic evidence hash and checks against capability baseline catalog; 5.~Scanner determines outcome: \texttt{passed} (zero findings/capabilities), \texttt{review-required} (capabilities requiring human attestation), or \texttt{needs-fixes} (blocking findings); 6.~Scanner persists immutable baseline scan record and posts diagnostic comment to issue. \\
\textbf{Alternate / Exception Flows} & \textbf{1a.~Repository Unreachable:} Scanner records transport error, retries with exponential backoff, and reports diagnostic failure if persistent. \textbf{3a.~High-Risk Finding Detected:} Scanner assigns \texttt{needs-fixes} status and blocks publication labeling. \\
\textbf{Related Requirements} & FR-22, FR-23, NFR-06, NFR-08 \\
\end{xltabular}

\begin{xltabular}{\textwidth}{@{}p{0.26\textwidth}X@{}}
\caption{Use-case specification: UC-14 Review and Publish Template Listing}\label{tab:uc14}\\
\toprule
\textbf{Component} & \textbf{Description} \\
\midrule
\endfirsthead
\multicolumn{2}{l}{\small\textit{(Continued from previous page)}}\\
\toprule
\textbf{Component} & \textbf{Description} \\
\midrule
\endhead
\bottomrule
\endlastfoot
\textbf{Actor} & Marketplace Maintainer (Secondary: GitHub Actions Publish Pipeline) \\
\textbf{System} & Marketplace Curation Engine \& Registry Pipeline \\
\textbf{Goal (Objective)} & Audit security baseline reports, attest capability reviews, and publish approved templates to the live public catalog. \\
\textbf{Preconditions} & Template submission has executed security baseline scan and maintainer has authorized write permissions. \\
\textbf{Trigger} & Maintainer applies \texttt{approved-and-verified} label to the submission issue. \\
\textbf{Basic Flow (Main Scenario)} & 1.~Maintainer inspects bot-authored security baseline report, visual previews, and repository provenance; 2.~If capabilities exist, maintainer verifies adherence to safety guidelines and applies approval attestation; 3.~Maintainer applies \texttt{approved-and-verified} label to submission issue; 4.~Publish workflow resolves exact label transition event identity and checks two-second timestamp compatibility window; 5.~Workflow clones repository, generates normalized WebP thumbnails via \texttt{sharp}, and updates \texttt{registry.json} and \texttt{catalog.json}; 6.~Automated test suite runs on updated catalog; upon passing, workflow commits changes to \texttt{main} branch and deploys Pages artifact. \\
\textbf{Alternate / Exception Flows} & \textbf{4a.~Stale Approval Event:} Event timestamp is ambiguous or multiple transitions detected; workflow fails closed and leaves catalog unchanged. \textbf{6a.~Regression Test Failure:} Catalog build or schema test fails; publication aborts, alert is posted to issue, and no commit is pushed. \\
\textbf{Related Requirements} & FR-20, FR-23, NFR-09 \\
\end{xltabular}

\begin{xltabular}{\textwidth}{@{}p{0.26\textwidth}X@{}}
\caption{Use-case specification: UC-15 Invoke Litho Operations via MCP Server}\label{tab:uc15}\\
\toprule
\textbf{Component} & \textbf{Description} \\
\midrule
\endfirsthead
\multicolumn{2}{l}{\small\textit{(Continued from previous page)}}\\
\toprule
\textbf{Component} & \textbf{Description} \\
\midrule
\endhead
\bottomrule
\endlastfoot
\textbf{Actor} & AI Assistant (Secondary: Litho MCP Daemon) \\
\textbf{System} & Litho Model Context Protocol (MCP) Server \\
\textbf{Goal (Objective)} & Allow LLM agents (e.g., Claude, Cursor, Antigravity) to autonomously build documents, validate schemas, retrieve templates, and extract source code over JSON-RPC. \\
\textbf{Preconditions} & Litho MCP server is running and connected via stdio or stream transport to an MCP-compliant host. \\
\textbf{Trigger} & AI assistant issues an MCP tool call (e.g., \texttt{build\_document}, \texttt{validate\_template}, \texttt{extract\_source}). \\
\textbf{Basic Flow (Main Scenario)} & 1.~AI agent discovers available Litho tool definitions and parameter schemas via MCP protocol discovery; 2.~Agent invokes selected tool (e.g., \texttt{build\_document}) passing document path and template identifier; 3.~MCP server validates input arguments against tool schema; 4.~Server dispatches request to underlying Litho CLI/compiler pipeline; 5.~Compilation succeeds and produces PDF artifact; 6.~MCP server serializes execution status, file path, and page metrics into JSON-RPC response delivered back to AI agent. \\
\textbf{Alternate / Exception Flows} & \textbf{3a.~Invalid Tool Arguments:} Server immediately returns JSON-RPC protocol error indicating missing or invalid parameters. \textbf{4a.~Compilation Failure:} Server catches pipeline error, extracts formatted error log, and returns diagnostic content to agent for conversational correction. \\
\textbf{Related Requirements} & FR-35--FR-41 \\
\end{xltabular}

\begin{xltabular}{\textwidth}{@{}p{0.26\textwidth}X@{}}
\caption{Use-case specification: UC-16 Automated PDF Export via Webhook}\label{tab:uc16}\\
\toprule
\textbf{Component} & \textbf{Description} \\
\midrule
\endfirsthead
\multicolumn{2}{l}{\small\textit{(Continued from previous page)}}\\
\toprule
\textbf{Component} & \textbf{Description} \\
\midrule
\endhead
\bottomrule
\endlastfoot
\textbf{Actor} & CMS Site / External Client (Secondary: Litho Conversion API Worker) \\
\textbf{System} & Litho Hosted Conversion API \& Webhook Service \\
\textbf{Goal (Objective)} & Trigger automated, headless document compilation from a content management system or Git commit webhook, receiving compiled PDF upon completion. \\
\textbf{Preconditions} & Client has valid API credentials and registered target webhook URL. \\
\textbf{Trigger} & CMS publishes new article or Git repository triggers push webhook. \\
\textbf{Basic Flow (Main Scenario)} & 1.~CMS delivers HTTP POST webhook request containing Markdown payload, template ID, and callback URL; 2.~Conversion API validates HMAC signature and IP rate limits; 3.~API queues asynchronous compilation job and returns HTTP 202 Accepted; 4.~Worker pulls job, compiles PDF in sandboxed execution container, and embeds source metadata; 5.~Worker uploads PDF to temporary signed storage; 6.~Worker dispatches HTTP POST callback to client's callback URL containing download link and build metadata. \\
\textbf{Alternate / Exception Flows} & \textbf{2a.~HMAC Signature Failure:} API rejects request immediately with HTTP 401 Unauthorized. \textbf{6a.~Callback Unreachable:} API retries callback with exponential backoff before logging delivery failure. \\
\textbf{Related Requirements} & FR-42, NFR-06, NFR-07 \\
\end{xltabular}

\begin{xltabular}{\textwidth}{@{}p{0.26\textwidth}X@{}}
\caption{Use-case specification: UC-17 Restore Source via Hosted API}\label{tab:uc17}\\
\toprule
\textbf{Component} & \textbf{Description} \\
\midrule
\endfirsthead
\multicolumn{2}{l}{\small\textit{(Continued from previous page)}}\\
\toprule
\textbf{Component} & \textbf{Description} \\
\midrule
\endhead
\bottomrule
\endlastfoot
\textbf{Actor} & CMS Site / External Client \\
\textbf{System} & Litho Hosted Conversion API \& Metadata Service \\
\textbf{Goal (Objective)} & Programmatically recover original Markdown source and template version from a Litho-generated PDF via HTTP REST endpoint. \\
\textbf{Preconditions} & Client has access to target PDF file and network connection to Litho API. \\
\textbf{Trigger} & Client issues HTTP POST request to \texttt{/api/v1/extract} with binary PDF payload. \\
\textbf{Basic Flow (Main Scenario)} & 1.~Client submits multi-part form data with PDF file to \texttt{/api/v1/extract}; 2.~API inspects PDF file header and XMP metadata streams; 3.~API locates embedded Litho source payload and verifies integrity hash; 4.~API decompresses raw Markdown content; 5.~API returns HTTP 200 OK with JSON response containing extracted Markdown source string, template ID, and template version. \\
\textbf{Alternate / Exception Flows} & \textbf{2a.~Non-Litho PDF:} API detects absence of Litho metadata; returns HTTP 422 Unprocessable Entity with error description. \textbf{3a.~Corrupted Stream:} API fails decompression or hash verification; returns HTTP 400 Bad Request. \\
\textbf{Related Requirements} & FR-43, NFR-05 \\
\end{xltabular}


\subsection{Activity Model}
\begin{guidance}
Model an important workflow using actions, control flows, decisions with guards,
parallel paths where required, and clear initial and final nodes. Use swimlanes when
responsibility changes are important.
\end{guidance}

The activity models below cover the authoring, publishing, marketplace, conversion, AI, MCP, and documentation workflows.

\diagram[dy]{Figures/activity/01-initialize-project}{ACT-01: Initialize a Litho project}
\diagram[dy]{Figures/activity/02-vscode-authoring}{ACT-02: Author a document using VS Code}
\diagram[dy]{Figures/activity/03-parse-validate}{ACT-03: Tokenize, parse, and validate Litho Markdown}
\diagram[dy]{Figures/activity/04-build-pdf}{ACT-04: Build a Markdown document into a PDF}
\diagram[dy]{Figures/activity/05-watch-rebuild}{ACT-05: Watch and rebuild a document}
\diagram[dy]{Figures/activity/06-extract-source}{ACT-06: Extract Markdown source from a PDF}
\diagram[dy]{Figures/activity/07-template-package}{ACT-07: Create and validate a template package}
\diagram[dy]{Figures/activity/08-publish-template}{ACT-08: Submit and publish a template through GitHub}
\diagram[dy]{Figures/activity/09-marketplace-template}{ACT-09: Browse, inspect, download, and use a marketplace template}
\diagram[dy]{Figures/activity/10-web-conversion}{ACT-10: Convert Markdown in the web application}
\diagram[dy]{Figures/activity/11-ai-authoring}{ACT-11: Use AI-assisted authoring}
\diagram[dy]{Figures/activity/12-security-review}{ACT-12: Perform AI-assisted template security review}
\diagram[dy]{Figures/activity/13-mcp-operation}{ACT-13: Invoke Litho through the MCP server}
\diagram[dy]{Figures/activity/14-documentation}{ACT-14: Build and publish documentation}

\subsection{State-Machine Model \applicability{State-dependent systems and games}}
\begin{guidance}
Model the lifecycle of one important stateful entity. Show states, triggering
events, guards where needed, actions, and initial and final states.
\end{guidance}
\diagram[dy]{Figures/state-machine/STATE-01-template-lifecycle}{STATE-01: Template package lifecycle}
\diagram[dy]{Figures/state-machine/STATE-02-compilation-job}{STATE-02: Document compilation job lifecycle}
\diagram[dy]{Figures/state-machine/STATE-03-web-conversion-session}{STATE-03: Web conversion workspace session}
\diagram[dy]{Figures/state-machine/STATE-04-submission-issue}{STATE-04: Template marketplace submission issue lifecycle}
\diagram[dy]{Figures/state-machine/STATE-05-author-workspace}{STATE-05: Author local workspace lifecycle}
\diagram[dy]{Figures/state-machine/STATE-06-mcp-session}{STATE-06: MCP server session lifecycle}
\diagram[dy]{Figures/state-machine/state-user-author}{Author user state model}
\diagram[dy]{Figures/state-machine/state-user-admin}{Administrator user state model}
\diagram[dy]{Figures/state-machine/state-user-publisher}{Publisher user state model}

\subsection{Sequence Model}
\begin{guidance}
Show how an important scenario is realized through time-ordered messages among
actors, system components, and external services. Use activation bars and combined
fragments such as \texttt{alt}, \texttt{opt}, or \texttt{loop} where appropriate.
\end{guidance}

The sequence models below show how parsing, template loading/binding, LaTeX generation, Tectonic compilation, and metadata embedding are realized, followed by catalog browsing, engagement recording, submissions, administration, and GitHub-based automation.

\diagram[dy]{Figures/sequence/01-browse-catalog}{Browse the public template catalog}
\diagram[dy]{Figures/sequence/02-search-filter}{Search and filter the catalog}
\diagram[dy]{Figures/sequence/03-template-detail}{Open a template detail page}
\diagram[dy]{Figures/sequence/04-record-engagement}{Record a view or download event}
\diagram[dy]{Figures/sequence/05-submit-template}{Submit a template from the public web form}
\diagram[dy]{Figures/sequence/06-admin-login}{Admin login}
\diagram[dy]{Figures/sequence/07-admin-create-template}{Admin creates a template}
\diagram[dy]{Figures/sequence/08-admin-approve-submission}{Admin approves a pending submission}
\diagram[dy]{Figures/sequence/09-admin-manage-taxonomy}{Admin manages taxonomy}
\diagram[dy]{Figures/sequence/10-admin-invite-admin}{Admin invites another administrator}
\diagram[dy]{Figures/sequence/11-admin-generate-registry}{Admin regenerates the registry snapshot}
\diagram[dy]{Figures/sequence/12-gh-route-issue}{Route a new plugin issue}
\diagram[dy]{Figures/sequence/13-gh-validate-submission}{Validate a plugin submission}
\diagram[dy]{Figures/sequence/14-gh-approve-publish}{Approve and publish a submission}
\diagram[dy]{Figures/sequence/15-gh-verify-plugin}{Verify a listed plugin}
\diagram[dy]{Figures/sequence/16-gh-refresh-catalog}{Refresh the public catalog}
\diagram[dy]{Figures/sequence/17-gh-delist-plugin}{Delist a plugin}
\diagram[dy]{Figures/sequence/18-deploy-github-pages}{Deploy the static site to GitHub Pages}
\diagram[dy]{Figures/sequence/19-mcp-tool-invocation}{AI assistant invokes Litho through the MCP server}

\section{Structural Design \applicability{Development, Hybrid, Game; Research if applicable}}
\begin{guidance}
Describe the static structure needed to realize the requirements. Keep models at a
level appropriate to the project and avoid including generated framework classes
that do not communicate a design decision.
\end{guidance}

\subsection{Class or Domain Model}
\begin{guidance}
Show important domain concepts or software classes, relevant attributes and
operations, relationships, multiplicities, and inheritance where justified.
\end{guidance}
\diagram[dy]{class-domain}{Class or domain diagram}

\subsection{Component Model}
\begin{guidance}
Show deployable or replaceable software components, their provided and required
interfaces, and dependencies.
\end{guidance}
\diagramplaceholder{Component diagram}{Show deployable or replaceable software components, their interfaces, and dependencies.}

\subsection{Package Model \applicability{If required}}
\begin{guidance}
Use a package diagram when the organization and dependencies of modules or layers
need clarification. Omit it when the component or class model already communicates
the structure adequately.
\end{guidance}
\diagram[dy]{package}{Package diagram}

\subsection{Deployment Model \applicability{Deployed, distributed, networked, or multi-device systems}}
\begin{guidance}
Show execution nodes, devices, environments, deployed artifacts, and communication
paths. Distinguish physical or virtual nodes from software components.
\end{guidance}
\diagram[dy]{deployment}{Deployment diagram}

\section{Data Design and Management}
\label{sec:data-design}

\subsection{Registry and Persistence Design \applicability{If applicable}}
\begin{guidance}
Describe persistent entities, identifiers, relationships, integrity constraints,
and important indexing or retention decisions. Use an ER diagram, schema, or other
notation appropriate to the selected database model.
\end{guidance}
Litho does not use an application database for the template registry. The published registry is derived from GitHub issues that maintainers review and label as listed, while each template's linked repository and releases hold its manifest, assets, and version history. A small application database stores only the per-template heart and download counters, keyed by template identifier. The diagram below models the resulting registry and counter records.
\diagramplaceholder{Data model or schema}{Show persistent entities, identifiers, relationships, integrity constraints, and important indexing or retention decisions.}

\section{Algorithm or Model Design}
\begin{guidance}
Describe only algorithms or models central to the project's contribution. Specify
inputs, outputs, principal processing steps, assumptions, training or configuration,
and computational or fairness considerations where relevant. Cite established
algorithms rather than presenting them as original work.
\end{guidance}

\begin{xltabular}{\textwidth}{@{}p{0.18\textwidth}X@{}}
\caption{Core algorithm or model}\label{tab:algorithm-model}\\
\toprule
\textbf{Item} & \textbf{Description} \\
\midrule
\endfirsthead
\multicolumn{2}{l}{\small\textit{(Continued from previous page)}}\\
\toprule
\textbf{Item} & \textbf{Description} \\
\midrule
\endhead
\multicolumn{2}{r}{\small\textit{(Continued on next page)}}\\
\endfoot
\bottomrule
\endlastfoot
Name and purpose & Template-binding Markdown-to-LaTeX compilation (core interpreter pipeline). \\
Inputs and outputs & Input: Litho Markdown (\texttt{.md}) and template ID/schema; Output: compilable \texttt{.tex} and compiled PDF with embedded metadata. \\
Processing steps & Tokenize → Parse to AST → Load template schema/manifest/assets → Validate/bind → Generate LaTeX → Compile via Tectonic → Embed source/metadata → Output PDF. \\
Assumptions and limits & Assumes CommonMark-compatible input with template-declared extensions; limited by Tectonic package coverage and template sandboxing. \\
Related requirements & FR-01--FR-08, NFR-04, NFR-05, NFR-06 \\
\end{xltabular}


\section{User Interface and Accessibility Design \applicability{User-facing Development, Hybrid, and Game projects; Research prototypes if applicable}}

\diagram[dy]{Figures/userflow/11-end-to-end-user-journey}{End-to-end user journey}
\diagram[dy]{Figures/userflow/01-user-md-to-pdf}{Markdown-to-PDF conversion flow}
\diagram[dy]{Figures/userflow/02-user-handles-conversion-error}{Conversion error handling flow}
\diagram[dy]{Figures/userflow/03-user-extracts-md-from-pdf}{Source extraction flow}
\diagram[dy]{Figures/userflow/04-user-browses-template}{Template marketplace browsing flow}
\diagram[dy]{Figures/userflow/05-litho-extension-working}{Litho extension workflow}
\diagram[dy]{Figures/userflow/06-cli-build-process}{CLI build process}
\diagram[dy]{Figures/userflow/07-template-creation-validation}{Template creation and validation flow}
\diagram[dy]{Figures/userflow/08-user-template-submission-review}{Template submission and review flow}
\diagram[dy]{Figures/userflow/09-ai-template-generation}{AI template generation flow}
\diagram[dy]{Figures/userflow/10-mcp-api-conversion}{MCP and API conversion flow}
\diagram[dy]{Figures/userflow/12-ai-security-review}{AI security review flow}

\diagram[dy]{Figures/wireframes/01-marketplace-catalogue}{Marketplace catalogue wireframe}
\diagram[dy]{Figures/wireframes/02-marketplace-template-detail}{Marketplace template detail wireframe}
\diagram[dy]{Figures/wireframes/03-marketplace-template-submission}{Marketplace template submission wireframe}
\diagram[dy]{Figures/wireframes/04-editor-start}{Editor start wireframe}
\diagram[dy]{Figures/wireframes/05-editor-workspace}{Editor workspace wireframe}
\diagram[dy]{Figures/wireframes/06-editor-pdf-export}{Editor PDF export wireframe}
\diagram[dy]{Figures/wireframes/07-documentation-home}{Documentation home wireframe}
\diagram[dy]{Figures/wireframes/08-documentation-cli}{Documentation CLI wireframe}
\diagram[dy]{Figures/wireframes/09-documentation-template-authoring}{Documentation template authoring wireframe}

Litho exposes three separate web applications: the template marketplace for browsing and discovering listed templates, the documentation website for installation and configuration guidance, and the conversion web app where authors write or paste Markdown, optionally use AI-assisted template generation, specification-to-documentation conversion, or template modification, select a template, and convert in either direction. Agents can also leverage the MCP server with the \texttt{skills.md} playbook to perform the same workflows programmatically. The same conversion capability is also available through the VS Code extension, CLI, MCP server, and hosted API.

\section{Security, Privacy, and Safety Design \applicability{If applicable}}
\begin{guidance}
Explain the design controls used to satisfy the corresponding requirements in
Chapter~\ref{chap:requirements}. Address relevant threats, authentication,
authorization, data protection, logging, secure failure, privacy, and safety. Do
not make general claims such as ``the system is secure'' without specific controls.
\end{guidance}

\begin{xltabular}{\textwidth}{@{}p{0.16\textwidth}X X p{0.19\textwidth}@{}}
\caption{Security, privacy, and safety controls}\label{tab:security-controls}\\
\toprule
\textbf{Design ID} & \textbf{Risk or requirement} & \textbf{Design control} & \textbf{Related requirement(s)} \\
\midrule
\endfirsthead
\multicolumn{4}{l}{\small\textit{(Continued from previous page)}}\\
\toprule
\textbf{Design ID} & \textbf{Risk or requirement} & \textbf{Design control} & \textbf{Related requirement(s)} \\
\midrule
\endhead
\multicolumn{4}{r}{\small\textit{(Continued on next page)}}\\
\endfoot
\bottomrule
\endlastfoot
SEC-01 & Arbitrary code execution in TeX engine & Invoke Tectonic with shell escapes strictly disabled and filesystem access restricted to the isolated build workspace. & FR-05, NFR-06, SE-01 \\
SEC-02 & Malicious template repository injection & Automated static security baseline scanning of submitted repository trees and declarative manifest schema validation prior to listing. & FR-14, FR-22, NFR-08, SE-02 \\
SEC-05 & AI template creation security & Automated AI-driven security review of templates created via the Template Maker; detected vulnerabilities are automatically filed as GitHub issues. & FR-28, NFR-08, SE-03 \\
SEC-03 & Denial of service and metric abuse & IP-prefix sliding-window rate limiting on public conversion endpoints and edge counter increments without persisting client IP addresses. & FR-48--FR-51, NFR-07, DIR-05 \\
SEC-04 & Cross-site scripting (XSS) and input poisoning & Strict HTML escaping of author-supplied Markdown metadata and enforcement of HTTPS URL protocol whitelists for template assets. & FR-24, FR-26, NFR-06, PR-01 \\
\end{xltabular}


\section{Design Decisions and Rationale}
\begin{guidance}
Record significant choices, alternatives considered, and reasons for selection.
Focus on decisions that materially affect quality, cost, risk, maintainability, or
project feasibility. Do not repeat the tools list from Section~\ref{sec:tools}.
\end{guidance}

\begin{xltabular}{\textwidth}{@{}p{0.14\textwidth}X X X@{}}
\caption{Design decisions}\label{tab:design-decisions}\\
\toprule
\textbf{ID} & \textbf{Decision} & \textbf{Alternatives considered} & \textbf{Rationale and trade-off} \\
\midrule
\endfirsthead
\multicolumn{4}{l}{\small\textit{(Continued from previous page)}}\\
\toprule
\textbf{ID} & \textbf{Decision} & \textbf{Alternatives considered} & \textbf{Rationale and trade-off} \\
\midrule
\endhead
\multicolumn{4}{r}{\small\textit{(Continued on next page)}}\\
\endfoot
\bottomrule
\endlastfoot
DD-01 & Adopt Tectonic as the self-contained XeTeX compilation engine. & Full TeX Live distribution (~5 GB installation), MiKTeX, or pdfLaTeX via Pandoc. & Bundles XeTeX, fontconfig, and on-demand CTAN package retrieval into a single portable binary without requiring system-wide TeX installations. Trade-off: initial build requires internet connectivity to cache packages. \\
DD-02 & Git-backed template distribution indexed via reviewed GitHub issues. & Centralized relational database (PostgreSQL/Supabase) or cloud object storage (S3) with custom auth. & Leverages distributed version control, issue tracking, and commit signing for open community curation with zero user credential management. Trade-off: catalog updates depend on GitHub API rate limits. \\
DD-03 & Embed original Markdown source into PDF XMP document metadata streams. & Heuristic reverse decompilation (transpiling generated LaTeX back to Markdown via regex or AST invertors). & Guarantees byte-for-byte lossless round-trip recovery of comments, frontmatter, and formatting. Trade-off: slightly increases compiled PDF file size (typically by 5--30 KB). \\
DD-04 & Jamstack static site generation with edge engagement counters (Next.js + Cloudflare D1). & Monolithic full-stack server-side rendered application (Node.js/Express or Django). & Maximizes global CDN delivery speed and eliminates server maintenance vulnerability surfaces, while edge counters handle lightweight reactions. Trade-off: new listings require catalog rebuild. \\
\end{xltabular}


% =============================================================================
